% META: {"id": "TEXP-ARI-CR-010", "chapitre": "TEXP-ARITHMETIQUE", "type_objet": "cours", "capacites_codes": ["C1", "C9"], "sources_inspiration": [], "mode_creation": "ex_nihilo", "status": "generated"}

\section{Divisibilité et PGCD}

% ============================================================
% STRATE 1 — Divisibilite, diviseurs
% ============================================================

\definition[1]{
Un entier $a\in\mathbb{Z}$ est \textbf{divisible} par $b\in\mathbb{Z}^*$ (ou $b$ \textbf{divise} $a$, noté $b\mid a$) s'il existe $k\in\mathbb{Z}$ tel que $a=bk$.
}

\exemple{
Les diviseurs positifs de $36$ sont $1,2,3,4,6,9,12,18,36$.
}

% BEGIN-VERIFY
% diviseurs_36 = [d for d in range(1, 37) if 36 % d == 0]
% assert diviseurs_36 == [1, 2, 3, 4, 6, 9, 12, 18, 36]
% END-VERIFY

% ============================================================
% STRATE 2 — PGCD et algorithme d'Euclide
% ============================================================

\definition[2]{
Le \textbf{PGCD} (plus grand commun diviseur) de deux entiers non tous deux nuls $a,b$ est le plus grand entier qui divise à la fois $a$ et $b$.
}

\theoreme[Algorithme d'Euclide]{
Pour $a\in\mathbb{Z}$ et $b\in\mathbb{N}^*$, en notant $r$ le reste de la division euclidienne de $a$ par $b$, on a $0\leq r<b$ et
$\mathrm{PGCD}(a,b)=\mathrm{PGCD}(b,r)$.
Les diviseurs communs sont conservés car $a=bq+r$.
En répétant le procédé, les restes non nuls décroissent strictement : l'algorithme termine.
À la dernière division, $r=0$ et le dernier diviseur $b$ est le PGCD.
}

\exemple{
Calculons $\mathrm{PGCD}(252,105)$ :
\[
  252 = 2\times105+42,\quad 105=2\times42+21,\quad 42=2\times21+0.
\]
Le dernier reste non nul est $21$ : $\mathrm{PGCD}(252,105)=21$.
}

Pour traiter aussi les entiers négatifs, le programme commence par les valeurs absolues : les diviseurs communs positifs ne changent pas. Si une entrée est nulle, l'autre étant non nulle, le PGCD est sa valeur absolue. Le couple $(0,0)$, exclu par la définition, est refusé.

\lstinputlisting[language=Python]{chapitres/TEXP-ARITHMETIQUE/python/euclide.py}

% BEGIN-VERIFY
% import hashlib
% import math
% from pathlib import Path
% import runpy
% depart = Path.cwd().resolve()
% relatif = "Mathematiques/manuel-maths/chapitres/TEXP-ARITHMETIQUE/python/euclide.py"
% fichier = next((p / relatif for p in (depart, *depart.parents) if (p / relatif).is_file()), None)
% assert fichier is not None, "source Python imprimee introuvable"
% assert hashlib.sha256(fichier.read_bytes()).hexdigest() == "0215fa7fcc2be276acaee67aeed83976e47db01d4ba33e7d9c2faf626fc82802"
% euclide = runpy.run_path(str(fichier))["euclide"]
% assert euclide(252, 105) == 21
% for a, b in ((252, 105), (17, 5), (-252, 105), (252, -105), (-252, -105), (0, -12), (-12, 0)):
%     assert euclide(a, b) == math.gcd(a, b)
% try:
%     euclide(0, 0)
% except ValueError:
%     pass
% else:
%     raise AssertionError("le couple (0, 0) doit etre refuse")
% END-VERIFY

\propriete[Couples d'entiers premiers entre eux]{
Deux entiers $a,b$ sont \textbf{premiers entre eux} si $\mathrm{PGCD}(a,b)=1$ (leur seul diviseur commun positif est $1$).
}

\erreurFrequente{
\textbf{Confondre \og premiers entre eux \fg{} et \og nombres premiers \fg{}.} Deux nombres peuvent être premiers entre eux sans être eux-mêmes des nombres premiers : par exemple, $8$ et $15$ sont premiers entre eux ($\mathrm{PGCD}(8,15)=1$), bien que ni $8$ ni $15$ ne soient des nombres premiers.
}
